% problem_id: S001-theorem-lazard
% source_id: S001 (Stacks Project)
% source_file: algebra.tex
% statement_locator: algebra.tex lines 19675--19679
% proof_locator: algebra.tex lines 19681--19740
% env: theorem
% id_note: label 在本来源内唯一
% status: complete (statement + author proof verbatim + reason + locators)
%
\begin{theorem}[Lazard's theorem]
\label{theorem-lazard}
Let $M$ be an $R$-module.  Then $M$ is flat if and only if it is the colimit of
a directed system of free finite $R$-modules.
\end{theorem}

\begin{proof}
A colimit of a directed system of flat modules is flat, as taking directed
colimits is exact and commutes with tensor product. Hence if $M$ is the colimit
of a directed system of free finite modules then $M$ is flat.

\medskip\noindent
For the converse, first recall that any module $M$ can be written as the
colimit of a directed system of finitely presented modules, in the following
way.  Choose a surjection $f: R^I \to M$ for some set $I$, and let $K$
be the kernel.  Let $E$ be the set of ordered pairs $(J, N)$ where $J$ is a
finite subset of $I$ and $N$ is a finitely generated submodule of $R^J \cap K$.
Then $E$ is made into a directed partially ordered set by defining $(J, N) \leq
(J', N')$ if and only if $J \subset J'$ and $N \subset N'$.  Define $M_e =
R^J/N$ for $e = (J, N)$, and define $f_{ee'}: M_e \to M_{e'}$ to be
the natural map for $e \leq e'$.  Then $(M_e, f_{ee'})$ is a directed system and
the natural maps $f_e: M_e \to M$ induce an isomorphism
$\colim_{e
\in E} M_e \xrightarrow{\cong} M$.

\medskip\noindent
Now suppose $M$ is flat.  Let $I = M \times \mathbf{Z}$, write $(x_i)$ for
the canonical basis of $R^{I}$, and take in the above discussion $f: R^I
\to M$ to be the map sending $x_i$ to the projection of $i$ onto $M$.
To prove the theorem it suffices to show that the $e \in E$ such that $M_e$
is free form a cofinal subset of $E$.  So let $e = (J, N) \in E$ be arbitrary.
By Lemma \ref{lemma-flat-factors-fp} there is a free finite module $F$ and maps
$h: R^J/N \to F$ and $g: F \to M$ such that the natural map
$f_e: R^J/N \to M$ factors as $R^J/N \xrightarrow{h} F
\xrightarrow{g} M$.  We are going to realize $F$ as $M_{e'}$ for some $e' \geq
e$.

\medskip\noindent
Let $\{ b_1, \ldots, b_n \}$ be a finite basis of $F$.  Choose $n$ distinct
elements $i_1, \ldots, i_n \in I$ such that $i_{\ell} \notin J$ for all $\ell$,
and such that the image of $x_{i_{\ell}}$ under $f: R^I \to M$ equals
the image of $b_{\ell}$ under $g: F \to M$.  This is possible since
every element of $M$ can be written as $f(x_i)$ for infinitely many
distinct $i \in I$ (by our choice of $I$).  Now let
$J' = J \cup \{i_1, \ldots , i_n \}$, and define $R^{J'}
\to F$ by $x_i \mapsto h(x_i)$ for $i \in J$ and $x_{i_{\ell}} \mapsto
b_{\ell}$ for $\ell = 1, \ldots, n$.  Let $N' = \Ker(R^{J'} \to F)$.
Observe:
\begin{enumerate}
\item The square
$$
\xymatrix{
R^{J'} \ar[r] \ar@{^{(}->}[d] & F \ar[d]^{g} \\
R^{I} \ar[r]_{f} & M
}
$$
is commutative,
%$R^{J'} \to F$ factors $f: R^I \to M$,
hence $N' \subset K = \Ker(f)$;
\item $R^{J'} \to F$ is a surjection onto a free finite module, hence
it splits and so $N'$ is finitely generated;
\item $J \subset J'$ and $N \subset N'$.
\end{enumerate}
By (1) and (2) $e' = (J', N')$ is in $E$, by (3) $e' \geq e$, and by
construction $M_{e'} = R^{J'}/N' \cong F$ is free.
\end{proof}
